\documentclass[../../main.tex]{subfiles}
\graphicspath{{\subfix{../../image/}}}

\begin{document}

\section{约束}\label{section:ch08-4-约束}

本节考察变分法在某些约束极小化问题上的应用，并特别讨论相应 Euler--Lagrange PDE 中 Lagrange 乘子的作用。

\subsection{非线性特征值问题}\label{subsection:ch08-4-1-非线性特征值}

首先研究带积分约束的问题。具体地，考察极小化能量泛函
\begin{align}
I[w]\triangleq\frac12\int_U|Dw|^2\,\mathrm{d}x\label{align:ch08-4-eq01}
\end{align}
（在满足 $w=0$ 在 $\partial U$ 上的所有函数 $w$ 中）的问题，但现还受附加边条件
\begin{align}
J[w]\triangleq\int_U G(w)\,\mathrm{d}x=0,\label{align:ch08-4-eq02}
\end{align}
其中 $G:\mathbb{R}\to\mathbb{R}$ 为给定光滑函数。今后写 $g=G'$。现假设
\begin{align}
|g(z)|\leqslant C(|z|+1),\label{align:ch08-4-eq03}
\end{align}
从而
\begin{align}
|G(z)|\leqslant C(|z|^2+1)\qquad (z\in\mathbb{R})\label{align:ch08-4-eq04}
\end{align}
对某常数 $C$。

还引入适当容许类
\begin{align*}
\mathcal{A}\triangleq\{w\in H_0^1(U)\mid J[w]=0\}.
\end{align*}
另设开集 $U$ 有界、连通且边界光滑。

\begin{theorem}[带约束极小元的存在性]\label{theorem:ch08-4-1-约束存在}
设容许集 $\mathcal{A}$ 非空。则存在 $u\in\mathcal{A}$ 满足
\begin{align*}
I[u]=\min_{w\in\mathcal{A}}I[w].
\end{align*}
\begin{proof}
如常选取极小化序列 $\{u_k\}_{k=1}^\infty\subset\mathcal{A}$，使
\begin{align*}
I[u_k]\to m=\inf_{w\in\mathcal{A}}I[w].
\end{align*}
于是如上可提取子列
\begin{align}
u_{k_j}\rightharpoonup u\quad\text{弱收敛于 }H_0^1(U),\label{align:ch08-4-eq05}
\end{align}
且 $I[u]\leqslant m$。一旦证明
\begin{align}
J[u]=0,\label{align:ch08-4-eq06}
\end{align}
从而 $u\in\mathcal{A}$，证毕。利用 \S5.7 的紧性理论，由\eqref{align:ch08-4-eq05} 推得
\begin{align}
u_{k_j}\to u\quad\text{在 }L^2(U)\text{ 中}.\label{align:ch08-4-eq07}
\end{align}
因此
\begin{align}
|J(u)|=|J(u)-J(u_k)|&\leqslant\int_U|G(u)-G(u_k)|\,\mathrm{d}x\notag\\
&\leqslant C\int_U|u-u_k|(1+|u|+|u_k|)\,\mathrm{d}x\quad\text{由 }\eqref{align:ch08-4-eq03}\notag\\
&\to 0\quad\text{当 }k\to\infty.\label{align:ch08-4-eq08}
\end{align}
\end{proof}
\end{theorem}

远比带约束极小元的单纯存在性更有趣的是考察相应的 Euler--Lagrange 方程。

\begin{theorem}[Lagrange 乘子]\label{theorem:ch08-4-2-Lagrange乘子}
设 $u\in\mathcal{A}$ 满足
\begin{align}
I[u]=\min_{w\in\mathcal{A}}I[w].\label{align:ch08-4-eq09}
\end{align}
则存在实数 $\lambda$ 使
\begin{align}
\int_U Du\cdot Dv\,\mathrm{d}x=\lambda\int_U g(u)v\,\mathrm{d}x\label{align:ch08-4-eq10}
\end{align}
对所有 $v\in H_0^1(U)$ 成立。

\begin{remark}\label{remark:ch08-4-1-非线性特征值}
于是 $u$ 是非线性边值问题
\begin{align}
\begin{cases}
-\Delta u=\lambda g(u)&\text{在 }U\text{ 中},\\
u=0&\text{在 }\partial U\text{ 上}
\end{cases}\label{align:ch08-4-eq11}
\end{align}
的弱解，其中 $\lambda$ 是对应于积分约束
\begin{align}
J[u]=0\label{align:ch08-4-eq12}
\end{align}
的 Lagrange 乘子。形如\eqref{align:ch08-4-eq11}、关于未知量 $(u,\lambda)$ 且 $u\not\equiv 0$ 的问题是一个\textbf{非线性特征值问题}。
\end{remark}

\begin{proof}
1. 固定任一函数 $v\in H_0^1(U)$。首先假设
\begin{align}
g(u)\text{ 在 }U\text{ 内 a.e. 不等于零。}\label{align:ch08-4-eq13}
\end{align}
然后任取函数 $w\in H_0^1(U)$ 满足
\begin{align}
\int_U g(u)w\,\mathrm{d}x\neq 0;\label{align:ch08-4-eq14}
\end{align}
由\eqref{align:ch08-4-eq13} 这是可能的。现写
\begin{align}
j(\tau,\sigma)\triangleq J[u+\tau v+\sigma w]=\int_U G(u+\tau v+\sigma w)\,\mathrm{d}x\qquad (\tau,\sigma\in\mathbb{R}).\label{align:ch08-4-eq15}
\end{align}
显然
\begin{align}
j(0,0)=\int_U G(u)\,\mathrm{d}x=0.\label{align:ch08-4-eq16}
\end{align}
此外 $j$ 为 $C^1$，且
\begin{align}
\frac{\partial j}{\partial\tau}(\tau,\sigma)=\int_U g(u+\tau v+\sigma w)v\,\mathrm{d}x,\label{align:ch08-4-eq17}
\end{align}
\begin{align}
\frac{\partial j}{\partial\sigma}(\tau,\sigma)=\int_U g(u+\tau v+\sigma w)w\,\mathrm{d}x.\label{align:ch08-4-eq18}
\end{align}
因此\eqref{align:ch08-4-eq14} 蕴涵
\begin{align}
\frac{\partial j}{\partial\sigma}(0,0)\neq 0.\label{align:ch08-4-eq19}
\end{align}
由隐函数定理（\S C.7），存在 $C^1$ 函数 $\phi:\mathbb{R}\to\mathbb{R}$ 使
\begin{align}
\phi(0)=0\label{align:ch08-4-eq20}
\end{align}
且
\begin{align}
j(\tau,\phi(\tau))=0\label{align:ch08-4-eq21}
\end{align}
对所有充分小的 $\tau$（设 $|\tau|\leqslant\tau_0$）成立。求导得
\begin{align*}
\frac{\partial j}{\partial\tau}(\tau,\phi(\tau))+\frac{\partial j}{\partial\sigma}(\tau,\phi(\tau))\phi'(\tau)=0,
\end{align*}
由此\eqref{align:ch08-4-eq17}、\eqref{align:ch08-4-eq18} 给出
\begin{align}
\phi'(0)=-\frac{\int_U g(u)v\,\mathrm{d}x}{\int_U g(u)w\,\mathrm{d}x}.\label{align:ch08-4-eq22}
\end{align}
2. 现令
\begin{align*}
w(\tau)\triangleq\tau v+\phi(\tau)w\qquad (|\tau|\leqslant\tau_0)
\end{align*}
并写
\begin{align*}
i(\tau)\triangleq I[u+w(\tau)].
\end{align*}
由于\eqref{align:ch08-4-eq21} 蕴涵 $J[u+w(\tau)]=0$，知 $u+w(\tau)\in\mathcal{A}$。故 $C^1$ 函数 $i(\cdot)$ 在 $0$ 处取最小值。于是
\begin{align}
0=i'(0)&=\left.\int_U(Du+\tau Dv+\phi(\tau)Dw)\cdot(Dv+\phi'(\tau)Dw)\,\mathrm{d}x\right|_{\tau=0}\notag\\
&=\int_U Du\cdot(Dv+\phi'(0)Dw)\,\mathrm{d}x.\label{align:ch08-4-eq23}
\end{align}
回顾\eqref{align:ch08-4-eq22} 并定义
\begin{align*}
\lambda\triangleq\frac{\int_U Du\cdot Dw\,\mathrm{d}x}{\int_U g(u)w\,\mathrm{d}x},
\end{align*}
由\eqref{align:ch08-4-eq23} 推得所需等式
\begin{align*}
\int_U Du\cdot Dv\,\mathrm{d}x=\lambda\int_U g(u)v\,\mathrm{d}x
\end{align*}
对所有 $v\in H_0^1(U)$ 成立。
3. 现设代替\eqref{align:ch08-4-eq13} 的是
\begin{align*}
g(u)=0\quad\text{a.e. 在 }U\text{ 中}.
\end{align*}
以有界函数逼近 $g$，推得 $DG(u)=g(u)Du=0$ a.e.。因此由于 $U$ 连通，$G(u)$ a.e. 为常数。又因 $J[u]=\int_U G(u)\,\mathrm{d}x=0$，故 $G(u)=0$ a.e.。由于 $u=0$ 在迹意义下在 $\partial U$ 上，故 $G(0)=0$。
但随后 $u=0$ a.e.，否则 $I[u]>I[0]=0$。由于 $g(u)=0$ a.e.，此情形下恒等式\eqref{align:ch08-4-eq10} 对任意 $\lambda$ 平凡成立。
\end{proof}
\end{theorem}

\subsection{单侧约束，变分不等式}\label{subsection:ch08-4-2-单侧约束}

现研究对每个 $x\in U$ 处 $u(x)$ 值施加某些逐点单侧约束的变分问题。为确定起见，考察极小化能量泛函
\begin{align}
I[w]\triangleq\int_U\frac12|Dw|^2-fw\,\mathrm{d}x\label{align:ch08-4-eq24}
\end{align}
（在属于集合
\begin{align}
\mathcal{A}\triangleq\{w\in H_0^1(U)\mid w\geqslant h\text{ a.e. 在 }U\text{ 中}\}\label{align:ch08-4-eq25}
\end{align}
的所有函数 $w$ 中）的问题，其中 $h:\overline{U}\to\mathbb{R}$ 为给定光滑函数，称为\textbf{障碍}。凸容许集 $\mathcal{A}$ 由 $H_0^1(U)$ 中满足单侧（即单边）约束 $w\geqslant h$ 的那些函数组成。还设 $f$ 为给定光滑函数。

\begin{theorem}[极小元的存在性]\label{theorem:ch08-4-3-障碍存在唯一}
设容许集 $\mathcal{A}$ 非空。则存在唯一函数 $u\in\mathcal{A}$ 满足
\begin{align*}
I[u]=\min_{w\in\mathcal{A}}I[w].
\end{align*}
\begin{proof}
1. 极小元的存在性由前面讨论的一般思想很容易得到。只需明确注意：若 $\{u_{k_j}\}_{j=1}^\infty\subset\mathcal{A}$ 是极小化序列且 $u_{k_j}\rightharpoonup u$ 弱收敛于 $H_0^1(U)$，则由紧性知 $u_{k_j}\to u$ 强收敛于 $L^2(U)$。由于 $u_{k_j}\geqslant h$ a.e.，故 $u\geqslant h$ a.e.。因此 $u\in\mathcal{A}$。
2. 现证唯一性。设 $u$ 与 $\widetilde{u}\in\mathcal{A}$ 为两个极小元，$u\neq\widetilde{u}$。则 $w\triangleq\frac{u+\widetilde{u}}2\in\mathcal{A}$，且
\begin{align*}
I[w]&=\int_U\frac12\left|\frac{Du+D\widetilde{u}}2\right|^2-f\!\left(\frac{u+\widetilde{u}}2\right)\mathrm{d}x\\
&=\int_U\frac18(|Du|^2+2Du\cdot D\widetilde{u}+|D\widetilde{u}|^2)-f\!\left(\frac{u+\widetilde{u}}2\right)\mathrm{d}x.
\end{align*}
由于 $2a\cdot b=|a|^2+|b|^2-|a-b|^2$，有
\begin{align*}
I[w]&=\int_U\frac18(2|Du|^2+2|D\widetilde{u}|^2-|Du-D\widetilde{u}|^2)-f\!\left(\frac{u+\widetilde{u}}2\right)\mathrm{d}x\\
&<\frac12\int_U\frac12|Du|^2-fu\,\mathrm{d}x+\frac12\int_U\frac12|D\widetilde{u}|^2-f\widetilde{u}\,\mathrm{d}x\\
&=\frac12 I[u]+\frac12 I[\widetilde{u}],
\end{align*}
严格不等式成立是因为 $u\not\equiv\widetilde{u}$。这是矛盾，因为 $u$ 与 $\widetilde{u}$ 是极小元。
\end{proof}
\end{theorem}

现计算 Euler--Lagrange 方程的类似物，在当前情形结果是一个不等式。

\begin{theorem}[极小元的变分刻画]\label{theorem:ch08-4-4-变分不等式}
设 $u\in\mathcal{A}$ 是
\begin{align*}
I[u]=\min_{w\in\mathcal{A}}I[w]
\end{align*}
的唯一解。则
\begin{align}
\int_U Du\cdot D(w-u)\,\mathrm{d}x\geqslant\int_U f(w-u)\,\mathrm{d}x\label{align:ch08-4-eq26}
\end{align}
对所有 $w\in\mathcal{A}$ 成立。称\eqref{align:ch08-4-eq26} 为\textbf{变分不等式}。
注意变分不等式中并未显式提及约束的 Lagrange 乘子。
\begin{proof}
1. 固定任一 $w\in\mathcal{A}$。则对每个 $0\leqslant\tau\leqslant 1$，
\begin{align*}
u+\tau(w-u)=(1-\tau)u+\tau w\in\mathcal{A},
\end{align*}
因为 $\mathcal{A}$ 凸。故若令
\begin{align*}
i(\tau)\triangleq I[u+\tau(w-u)],
\end{align*}
则对所有 $0\leqslant\tau\leqslant 1$ 有 $i(0)\leqslant i(\tau)$。因此
\begin{align}
i'(0)\geqslant 0.\label{align:ch08-4-eq27}
\end{align}
2. 若 $0<\tau\leqslant 1$，
\begin{align*}
\frac{i(\tau)-i(0)}{\tau}&=\frac1\tau\int_U\frac{|Du+\tau D(w-u)|^2-|Du|^2}{2}-f(u+\tau(w-u)-u)\,\mathrm{d}x\\
&=\int_U Du\cdot D(w-u)+\frac{\tau|D(w-u)|^2}{2}-f(w-u)\,\mathrm{d}x.
\end{align*}
故\eqref{align:ch08-4-eq27} 蕴涵
\begin{align*}
0\leqslant i'(0)=\int_U Du\cdot D(w-u)-f(w-u)\,\mathrm{d}x.
\end{align*}
\end{proof}
\end{theorem}

注意由于实际上只能取"单侧"变分（远离约束），故得到不等式\eqref{align:ch08-4-eq27}。

\paragraph{变分不等式的解释。}为对变分不等式\eqref{align:ch08-4-eq26} 有所洞见，我们不加证明地引用一个正则性断言（见 Kinderlehrer--Stampacchia [K-S]）：若 $\partial U$ 光滑，则 $u\in W^{2,\infty}(U)$。于是集合
\begin{align*}
O\triangleq\{x\in U\mid u(x)>h(x)\}
\end{align*}
开，而
\begin{align*}
C\triangleq\{x\in U\mid u(x)=h(x)\}
\end{align*}
（相对）闭。

我们断言事实上 $u\in C^\infty(O)$ 且
\begin{align}
-\Delta u=f\quad\text{在 }O\text{ 中}.\label{align:ch08-4-eq28}
\end{align}
为此，固定任一检验函数 $v\in C_c^\infty(O)$。则若 $|\tau|$ 充分小，$w\triangleq u+\tau v\geqslant h$，故 $w\in\mathcal{A}$。于是\eqref{align:ch08-4-eq26} 蕴涵 $\tau\int_O Du\cdot Dv-fv\,\mathrm{d}x\geqslant 0$。此不等式对所有充分小的 $\tau$（无论正负）成立，故事实上
\begin{align*}
\int_O Du\cdot Dv-fv\,\mathrm{d}x=0
\end{align*}
对所有 $v\in C_c^\infty(O)$ 成立。因此 $u$ 是\eqref{align:ch08-4-eq28} 的弱解，线性正则性理论（\S6.3）表明 $u\in C^\infty(O)$。

现若 $v\in C_c^\infty(U)$ 满足 $v\geqslant 0$ 且 $0<\tau\leqslant 1$，则 $w\triangleq u+\tau v\in\mathcal{A}$，从而 $\int_U Du\cdot Dv-fv\,\mathrm{d}x\geqslant 0$。但由于 $u\in W^{2,\infty}(U)$，分部积分可推得对所有非负函数 $v\in C_c^\infty(U)$，$\int_U(-\Delta u-f)v\,\mathrm{d}x\geqslant 0$。故
\begin{align}
-\Delta u\geqslant f\quad\text{a.e. 在 }U\text{ 中}.\label{align:ch08-4-eq29}
\end{align}

\paragraph{障碍问题的自由边界。}由\eqref{align:ch08-4-eq28}、\eqref{align:ch08-4-eq29} 总结我们的结论：
\begin{align}
\begin{cases}
u\geqslant h,\ -\Delta u\geqslant f&\text{a.e. 在 }U\text{ 中},\\
-\Delta u=f&\text{在 }U\cap\{u>h\}\text{ 上}.
\end{cases}\label{align:ch08-4-eq30}
\end{align}
\paragraph{自由边界。}集合
\begin{align*}
F\triangleq\partial O\cap U
\end{align*}
称为\textbf{自由边界}。应用数学中许多有趣问题涉及带自由边界的偏微分方程。其中那些可改写成变分不等式形式的问题变得相对容易研究，尤其因为不等式\eqref{align:ch08-4-eq30} 中并未显式提及自由边界。应用出现在 Brown 运动的停时最优控制问题、地下水水文学、塑性理论等中。见 Kinderlehrer--Stampacchia [K-S]、Friedman [Fr3]。

\subsection{调和映射}\label{subsection:ch08-4-3-调和映射}

接着考察以映到球面的调和映射为例的逐点约束情形。现感兴趣的问题是在属于容许类
\begin{align}
\mathcal{A}\triangleq\{\mathbf{w}\in H^1(U;\mathbb{R}^m)\mid \mathbf{w}=\mathbf{g}\text{ 在 }\partial U\text{ 上},\ |\mathbf{w}|=1\text{ a.e.}\}\label{align:ch08-4-eq32}
\end{align}
的所有函数上极小化能量
\begin{align}
I[\mathbf{w}]\triangleq\frac12\int_U|D\mathbf{w}|^2\,\mathrm{d}x.\label{align:ch08-4-eq31}
\end{align}
其想法是：在从 $U\subset\mathbb{R}^n$ 到单位球面 $S^{m-1}=\partial B(0,1)\subset\mathbb{R}^m$ 的所有适当映射上极小化能量。这个问题及其变体作为液晶行为的简化模型出现。

\begin{figure}[htbp]
\centering
\includegraphics[width=0.6\textwidth]{ch08-fig-harmonic-map-sphere.png}
\caption{映到球面的调和映射}\label{figure:ch08-harmonic-map-sphere}
\end{figure}

容易验证，只要 $\mathcal{A}\neq\emptyset$，$\mathcal{A}$ 中至少存在一个极小元。

\begin{theorem}[调和映射的 Euler--Lagrange 方程]\label{theorem:ch08-4-5-调和映射EL}
设 $\mathbf{u}\in\mathcal{A}$ 满足
\begin{align*}
I[\mathbf{u}]=\min_{\mathbf{w}\in\mathcal{A}}I[\mathbf{w}].
\end{align*}
则
\begin{align}
\int_U D\mathbf{u}:D\mathbf{v}\,\mathrm{d}x=\int_U|D\mathbf{u}|^2\mathbf{u}\cdot\mathbf{v}\,\mathrm{d}x\label{align:ch08-4-eq33}
\end{align}
对每个 $\mathbf{v}\in H_0^1(U;\mathbb{R}^m)\cap L^\infty(U;\mathbb{R}^m)$ 成立。

把\eqref{align:ch08-4-eq33} 解释为 $\mathbf{u}=(u^1,\cdots,u^m)$ 是边值问题
\begin{align}
\begin{cases}
-\Delta\mathbf{u}=|D\mathbf{u}|^2\mathbf{u}&\text{在 }U\text{ 中},\\
\mathbf{u}=\mathbf{g}&\text{在 }\partial U\text{ 上}
\end{cases}\label{align:ch08-4-eq34}
\end{align}
的弱解。函数 $\lambda=|D\mathbf{u}|^2$ 是对应于逐点约束 $|\mathbf{u}|=1$ 的 Lagrange 乘子。请注意：对于单个积分约束（\S8.4.1），Lagrange 乘子是一个数；但对于逐点约束，它是一个函数。
\begin{proof}
1. 固定 $\mathbf{v}\in H_0^1(U;\mathbb{R}^m)\cap L^\infty(U;\mathbb{R}^m)$。由于 $|\mathbf{u}|=1$ a.e.，对每个充分小的 $\tau$，有 $|\mathbf{u}+\tau\mathbf{v}|\neq 0$ a.e.。因此
\begin{align}
\mathbf{v}(\tau)\triangleq\frac{\mathbf{u}+\tau\mathbf{v}}{|\mathbf{u}+\tau\mathbf{v}|}\in\mathcal{A}.\label{align:ch08-4-eq35}
\end{align}
于是
\begin{align*}
i(\tau)\triangleq I[\mathbf{v}(\tau)]
\end{align*}
在 $\tau=0$ 处取最小值，故如常
\begin{align}
i'(0)=0.\label{align:ch08-4-eq36}
\end{align}
2. 现
\begin{align}
i'(0)=\int_U D\mathbf{u}:D\mathbf{v}'(0)\,\mathrm{d}x\qquad ('=\frac{\mathrm{d}}{\mathrm{d}\tau}).\label{align:ch08-4-eq37}
\end{align}
但由\eqref{align:ch08-4-eq35} 直接算得
\begin{align*}
\mathbf{v}'(\tau)=\frac{\mathbf{v}}{|\mathbf{u}+\tau\mathbf{v}|}-\frac{[(\mathbf{u}+\tau\mathbf{v})\cdot\mathbf{v}](\mathbf{u}+\tau\mathbf{v})}{|\mathbf{u}+\tau\mathbf{v}|^3},
\end{align*}
故 $\mathbf{v}'(0)=\mathbf{v}-(\mathbf{u}\cdot\mathbf{v})\mathbf{u}$。把此等式代入\eqref{align:ch08-4-eq36}、\eqref{align:ch08-4-eq37}，得
\begin{align}
0=\int_U D\mathbf{u}:D\mathbf{v}-D\mathbf{u}:D((\mathbf{u}\cdot\mathbf{v})\mathbf{u})\,\mathrm{d}x.\label{align:ch08-4-eq38}
\end{align}
然而由于 $|\mathbf{u}|^2\equiv 1$，有
\begin{align*}
(D\mathbf{u})^T\mathbf{u}=0.
\end{align*}
利用此事实，可验证
\begin{align*}
D\mathbf{u}:D((\mathbf{u}\cdot\mathbf{v})\mathbf{u})=|D\mathbf{u}|^2(\mathbf{u}\cdot\mathbf{v})\quad\text{a.e. 在 }U\text{ 中}.
\end{align*}
把此恒等式用于\eqref{align:ch08-4-eq38} 即得\eqref{align:ch08-4-eq33}。
\end{proof}
\end{theorem}

\subsection{不可压缩性}\label{subsection:ch08-4-4-不可压缩性}

\subsubsection*{a. Stokes 问题}

设 $U\subset\mathbb{R}^3$ 开、有界、单连通，令
\begin{align*}
I[\mathbf{w}]\triangleq\int_U\frac12|D\mathbf{w}|^2-\mathbf{f}\cdot\mathbf{w}\,\mathrm{d}x,
\end{align*}
其中 $\mathbf{w}$ 属于
\begin{align*}
\mathcal{A}\triangleq\{\mathbf{w}\in H_0^1(U;\mathbb{R}^3)\mid \operatorname{div}\mathbf{w}=0\text{ 在 }U\text{ 中}\}.
\end{align*}
此处 $\mathbf{f}\in L^2(U;\mathbb{R}^3)$ 给定。

用通常方法证明存在唯一极小元 $\mathbf{u}\in\mathcal{A}$ 并无困难。我们把 $\mathbf{u}$ 解释为区域 $U$ 内受外力 $\mathbf{f}$ 作用的定常流体流动的速度场。约束 $\operatorname{div}\mathbf{u}=0$ 保证流动不可压缩；见本小节末的注记。

约束在 Euler--Lagrange 方程中如何体现？

\begin{theorem}[压力作为 Lagrange 乘子]\label{theorem:ch08-4-6-压力乘子}
存在标量函数 $p\in L^2_{\mathrm{loc}}(U)$ 使
\begin{align}
\int_U D\mathbf{u}:D\mathbf{v}\,\mathrm{d}x=\int_U p\operatorname{div}\mathbf{v}+\mathbf{f}\cdot\mathbf{v}\,\mathrm{d}x\label{align:ch08-4-eq39}
\end{align}
对所有在 $U$ 内具有紧支集的 $\mathbf{v}\in H^1(U;\mathbb{R}^3)$ 成立。

把\eqref{align:ch08-4-eq39} 解释为 $(\mathbf{u},p)$ 是 Stokes 问题
\begin{align}
\begin{cases}
-\Delta\mathbf{u}=\mathbf{f}-Dp&\text{在 }U\text{ 中},\\
\operatorname{div}\mathbf{u}=0&\text{在 }U\text{ 中},\\
\mathbf{u}=0&\text{在 }\partial U\text{ 上}
\end{cases}\label{align:ch08-4-eq40}
\end{align}
的弱解。函数 $p$ 是压力，作为对应于不可压缩条件 $\operatorname{div}\mathbf{u}=0$ 的 Lagrange 乘子出现。
\begin{proof}
1. 首先设 $\mathbf{v}\in\mathcal{A}$。则对每个 $\tau\in\mathbb{R}$，$\mathbf{u}+\tau\mathbf{v}\in\mathcal{A}$。故
\begin{align}
0=i'(0)=\int_U D\mathbf{u}:D\mathbf{v}-\mathbf{f}\cdot\mathbf{v}\,\mathrm{d}x.\label{align:ch08-4-eq41}
\end{align}
2. 现固定 $V\subset \subset U$，$V$ 光滑且单连通，选 $\mathbf{w}\in H_0^1(V;\mathbb{R}^3)$ 满足 $\operatorname{div}\mathbf{w}=0$。取 $0<\varepsilon<\operatorname{dist}(V,\partial U)$，在\eqref{align:ch08-4-eq41} 中令 $\mathbf{v}=\mathbf{v}_\varepsilon\triangleq\eta_\varepsilon*\mathbf{w}$，其中 $\eta_\varepsilon$ 为通常磨光子，$\mathbf{w}$ 在 $U\setminus V$ 中定义为零。则
\begin{align}
0=\int_U D\mathbf{u}:D\mathbf{v}_\varepsilon-\mathbf{f}\cdot\mathbf{v}_\varepsilon\,\mathrm{d}x=\int_U D\mathbf{u}_\varepsilon:D\mathbf{w}-\mathbf{f}_\varepsilon\cdot\mathbf{w}\,\mathrm{d}x,\label{align:ch08-4-eq42}
\end{align}
其中
\begin{align}
\mathbf{u}_\varepsilon\triangleq\eta_\varepsilon*\mathbf{u},\qquad \mathbf{f}_\varepsilon\triangleq\eta_\varepsilon*\mathbf{f}.\label{align:ch08-4-eq43}
\end{align}
由于 $\mathbf{u}_\varepsilon$ 光滑，\eqref{align:ch08-4-eq42} 蕴涵
\begin{align}
\int_V(-\Delta\mathbf{u}_\varepsilon-\mathbf{f}_\varepsilon)\cdot\mathbf{w}\,\mathrm{d}x=0\label{align:ch08-4-eq44}
\end{align}
对每个 $\mathbf{w}\in H_0^1(V;\mathbb{R}^3)$ 且 $\operatorname{div}\mathbf{w}=0$ 成立。
3. 固定任一光滑向量场 $\boldsymbol{\zeta}\in C_c^\infty(V;\mathbb{R}^3)$，在\eqref{align:ch08-4-eq44} 中令 $\mathbf{w}=\operatorname{curl}\boldsymbol{\zeta}$。这是合法的，因为 $\operatorname{div}\mathbf{w}=\operatorname{div}(\operatorname{curl}\boldsymbol{\zeta})=0$。临时写 $\mathbf{h}\triangleq\Delta\mathbf{u}_\varepsilon+\mathbf{f}_\varepsilon$，得
\begin{align*}
0=\int_V\mathbf{h}\cdot\operatorname{curl}\boldsymbol{\zeta}\,\mathrm{d}x=\int_V h^1(\zeta^3_{x_2}-\zeta^2_{x_3})+h^2(\zeta^1_{x_3}-\zeta^3_{x_1})+h^3(\zeta^2_{x_1}-\zeta^1_{x_2})\,\mathrm{d}x,
\end{align*}
其中 $\mathbf{h}=(h^1,h^2,h^3)$，$\boldsymbol{\zeta}=(\zeta^1,\zeta^2,\zeta^3)$。分部积分得
\begin{align*}
0=\int_V\zeta^1(h^3_{x_2}-h^2_{x_3})+\zeta^2(h^1_{x_3}-h^3_{x_1})+\zeta^3(h^2_{x_1}-h^1_{x_2})\,\mathrm{d}x.
\end{align*}
由于 $\zeta^1,\zeta^2,\zeta^3\in C_c^\infty(V)$ 任意，推得 $\operatorname{curl}\mathbf{h}=0$ 在 $V$ 中。由于 $V$ 单连通，故存在 $V$ 中光滑函数 $p_\varepsilon$ 使
\begin{align}
Dp_\varepsilon=\mathbf{h}=\Delta\mathbf{u}_\varepsilon+\mathbf{f}_\varepsilon\quad\text{在 }V\text{ 中}.\label{align:ch08-4-eq45}
\end{align}
4. 必要时可给 $p_\varepsilon$ 加常数使 $\int_V p_\varepsilon\,\mathrm{d}x=0$。鉴于此归一化，存在光滑向量场 $\mathbf{v}_\varepsilon:V\to\mathbb{R}^3$ 解
\begin{align}
\begin{cases}
\operatorname{div}\mathbf{v}_\varepsilon=p_\varepsilon&\text{在 }V\text{ 中},\\
\mathbf{v}_\varepsilon=0&\text{在 }\partial V\text{ 上}.
\end{cases}\label{align:ch08-4-eq46}
\end{align}
此外有估计
\begin{align}
\|\mathbf{v}_\varepsilon\|_{H^1(V;\mathbb{R}^3)}\leqslant C\|p_\varepsilon\|_{L^2(V)},\label{align:ch08-4-eq47}
\end{align}
常数 $C$ 仅依赖于 $V$。（我们略去向量场 $\mathbf{v}_\varepsilon$ 存在性的证明。其构造既繁琐又需要 Neumann 型边界条件下 Laplace 方程某些估计的知识，超出本书范围。细节见 Dacorogna--Moser (Ann.\ Inst.\ H Poincaré 7 (1990), 1--26)。）
现计算
\begin{align*}
\int_V(p_\varepsilon)^2\,\mathrm{d}x&=\int_V p_\varepsilon\operatorname{div}\mathbf{v}_\varepsilon\,\mathrm{d}x\quad\text{由 }\eqref{align:ch08-4-eq46}\\
&=-\int_V Dp_\varepsilon\cdot\mathbf{v}_\varepsilon\,\mathrm{d}x\\
&=\int_V(-\Delta\mathbf{u}_\varepsilon-\mathbf{f}_\varepsilon)\cdot\mathbf{v}_\varepsilon\,\mathrm{d}x\quad\text{由 }\eqref{align:ch08-4-eq45}\\
&=\int_V D\mathbf{u}_\varepsilon:D\mathbf{v}_\varepsilon-\mathbf{f}_\varepsilon\cdot\mathbf{v}_\varepsilon\,\mathrm{d}x\\
&\leqslant\|\mathbf{v}_\varepsilon\|_{H^1(V;\mathbb{R}^3)}(\|\mathbf{u}_\varepsilon\|_{H^1(V)}+\|\mathbf{f}_\varepsilon\|_{L^2(V)})\\
&\leqslant C\|p_\varepsilon\|_{L^2(V)}(\|\mathbf{u}\|_{H^1(U)}+\|\mathbf{f}\|_{L^2(U)})\quad\text{由 }\eqref{align:ch08-4-eq47}.
\end{align*}
故
\begin{align}
\|p_\varepsilon\|_{L^2(V)}\leqslant C(\|\mathbf{u}\|_{H^1(U)}+\|\mathbf{f}\|_{L^2(U)}).\label{align:ch08-4-eq48}
\end{align}
5. 由估计\eqref{align:ch08-4-eq48}，存在子列 $\varepsilon_j\to 0$ 使
\begin{align}
p_{\varepsilon_j}\rightharpoonup p\quad\text{弱收敛于 }L^2(V)\label{align:ch08-4-eq49}
\end{align}
某 $p\in L^2(V)$。由\eqref{align:ch08-4-eq45} 蕴涵
\begin{align*}
\int_V D\mathbf{u}_\varepsilon:D\mathbf{v}\,\mathrm{d}x=\int_V p_\varepsilon\operatorname{div}\mathbf{v}+\mathbf{f}_\varepsilon\cdot\mathbf{v}\,\mathrm{d}x
\end{align*}
对所有 $\mathbf{v}\in H_0^1(V;\mathbb{R}^3)$ 成立。令 $\varepsilon=\varepsilon_j\to 0$，得
\begin{align}
\int_V D\mathbf{u}:D\mathbf{v}\,\mathrm{d}x=\int_V p\operatorname{div}\mathbf{v}+\mathbf{f}\cdot\mathbf{v}\,\mathrm{d}x.\label{align:ch08-4-eq50}
\end{align}
6. 最后选取集合序列 $V_k\subset \subset U$（$k=1,\cdots$）如上，$V_1\subset V_2\subset V_3\subset\cdots$ 且 $U=\bigcup_{k=1}^\infty V_k$。利用第 2--5 步，找到 $p_k\in L^2(V_k)$（$k=1,\cdots$）使
\begin{align}
\int_{V_k}D\mathbf{u}:D\mathbf{v}\,\mathrm{d}x=\int_{V_k}p_k\operatorname{div}\mathbf{v}+\mathbf{f}\cdot\mathbf{v}\,\mathrm{d}x\label{align:ch08-4-eq51}
\end{align}
对每个 $\mathbf{v}\in H_0^1(V_k;\mathbb{R}^3)$ 成立。必要时给每个 $p_k$ 加常数，由\eqref{align:ch08-4-eq51} 推得：若 $1\leqslant l\leqslant k$，则 $p_k=p_l$ 在 $V_l$ 上。最后定义 $p=p_k$ 在 $V_k$ 上（$k=1,\cdots$）。
\end{proof}
\end{theorem}

\subsubsection*{b. 不可压缩非线性弹性}

现回到 \S8.2.4 中讨论过的非线性弹性模型。设 $\mathbf{u}$ 表示静止构型为 $U$ 的弹性体的位移。现设弹性体不可压缩，这意味着
\begin{align*}
\det D\mathbf{u}=1.
\end{align*}
因此设给定能量密度函数 $L:\mathbb{M}^{3\times 3}\times U\to\mathbb{R}$，并考虑在容许集
\begin{align*}
\mathcal{A}\triangleq\{\mathbf{w}\in W^{1,q}(U;\mathbb{R}^3)\mid \mathbf{w}=\mathbf{g}\text{ 在 }\partial U\text{ 上},\ \det D\mathbf{w}=1\text{ a.e.}\}
\end{align*}
（某 $q>3$）中的所有 $\mathbf{w}$ 上极小化弹性能
\begin{align*}
I[\mathbf{w}]\triangleq\int_U L(D\mathbf{w},x)\,\mathrm{d}x
\end{align*}
的问题。

\begin{theorem}[带行列式约束的极小元]\label{theorem:ch08-4-7-不可压缩极小元}
设映射 $P\mapsto L(P,x)$ 凸，且 $L$ 满足强制性条件
\begin{align*}
L(P,x)\geqslant\alpha|P|^q-\beta\qquad (P\in\mathbb{M}^{3\times 3},\ x\in U)
\end{align*}
某 $\alpha>0$，$\beta\geqslant 0$。最后设 $\mathcal{A}\neq\emptyset$。
则存在 $\mathbf{u}\in\mathcal{A}$ 满足
\begin{align*}
I[\mathbf{u}]=\min_{\mathbf{w}\in\mathcal{A}}I[\mathbf{w}].
\end{align*}
\begin{proof}
如常选取极小化序列，
\begin{align*}
\mathbf{u}_{k_j}\rightharpoonup\mathbf{u}\quad\text{弱收敛于 }W^{1,q}(U;\mathbb{R}^3).
\end{align*}
由于
\begin{align*}
I[\mathbf{u}]\leqslant\liminf_{j\to\infty}I[\mathbf{u}_{k_j}],
\end{align*}
只需证明 $\mathbf{u}\in\mathcal{A}$。然而由 \S8.2.4 的引理，$\det D\mathbf{u}_{k_j}\rightharpoonup\det D\mathbf{u}$ 弱收敛于 $L^{q/n}(U)$，故 $\det D\mathbf{u}=1$ a.e.，正如所需。
\end{proof}
\end{theorem}

\begin{remark}[不可压缩性注记]\label{remark:ch08-4-2-不可压缩}
Stokes 问题的不可压缩条件是
\begin{align}
\operatorname{div}\mathbf{u}=0,\label{align:ch08-4-eq52}
\end{align}
而非线性弹性问题的是
\begin{align}
\det D\mathbf{u}=1,\label{align:ch08-4-eq53}
\end{align}
这看起来可能奇怪。解释在于 $\mathbf{u}$ 在\eqref{align:ch08-4-eq52} 中表示速度，而在\eqref{align:ch08-4-eq53} 中表示位移。更一般地，若 $\mathbf{b}$ 是速度场（例如流体的），则通过解 ODE
\begin{align*}
\begin{cases}
\dot{\mathbf{x}}(t)=\mathbf{b}(\mathbf{x}(t),t)&(t\in\mathbb{R}),\\
\mathbf{x}(0)=x
\end{cases}
\end{align*}
计算初始位于点 $x$ 的粒子的运动。写 $\mathbf{x}(t)=\mathbf{x}(t,x)$ 以显式显示对初位置 $x$ 的依赖。则对每个 $t>0$，映射 $x\mapsto\mathbf{x}(t,x)$ 保体积当且仅当
\begin{align*}
J(x,t)=\det D_x\mathbf{x}(t,x)=1
\end{align*}
对所有 $x$。显然 $J(x,0)=1$，计算验证 Euler 公式：
\begin{align*}
J_t=(\operatorname{div}\mathbf{b})(\mathbf{x},t)J,
\end{align*}
散度关于空间变量取。故若 $\operatorname{div}\mathbf{b}\equiv 0$，流动保体积。
\end{remark}

\end{document}
